\section{问题三模型建立及求解}

\subsection{双星协同与初始轨道联合设计模型的建立}

问题二中，单颗观测星 R1 从目标队列中部出发，必须先向西扫过西侧目标、再折返向东扫过东侧目标，30 天内跨越约 $38^\circ$ 的相位弧段，导致折返和高速追赶消耗了大量燃料。问题三引入一颗提前入轨的高轨卫星 R2，与 R1 共同在 30 天内观测全部 30 颗目标。我们将 30 条时变可观测管道划分为两个互不相交的子集，由 R1 和 R2 分别执行单向巡游；同时把 R2 在初始历元 $t_0$ 时刻的轨道状态作为连续决策变量，与其 30 天内的全部脉冲机动联合优化。双星各自免去了折返过程，整体思路如图~\ref{fig:q3abs} 所示。

\begin{figure}[!htb]
  \centering
  \includegraphics[width=0.95\textwidth]{fig11_q3abs.pdf}
  \caption{问题三的抽象：目标经度分组、R2 初始轨道设计与双星单向并行巡游}
  \label{fig:q3abs}
\end{figure}

\subsubsection{决策变量定义}

双星协同方案由目标分配、各自的访问序列与机动序列，以及 R2 的初始轨道状态组成。记 R1 与 R2 的访问序列分别为 $\pi^{(1)} = (\pi^{(1)}_1, \dots, \pi^{(1)}_{K_1})$ 与 $\pi^{(2)} = (\pi^{(2)}_1, \dots, \pi^{(2)}_{K_2})$，两星各段的机动与观测变量仍按式~\eqref{eq:q2var} 记为 $\bm z^{(1)}_k$（$k = 1, \dots, K_1$）和 $\bm z^{(2)}_k$（$k = 1, \dots, K_2$）。由于 R2 提前发射入轨，其在 $t_0$ 时刻的六维位置速度状态
\begin{equation}
  \bm x_0^{(2)} = \big(\bm r_0^{(2)},\ \bm v_0^{(2)}\big) \in \mathbb R^6
  \label{eq:q3x0}
\end{equation}
也是决策变量。完整方案记为 $\mathcal P = \big(\pi^{(1)}, \bm z^{(1)};\ \bm x_0^{(2)}, \pi^{(2)}, \bm z^{(2)}\big)$。

\subsubsection{目标函数设置}

按照第 4.4 节的两阶段字典序目标，第一阶段最大化两星合计观测到的不同目标数 $K = K_1 + K_2$，第二阶段在观测数最大的方案中最小化两星累计速度增量之和
\begin{equation}
  J(\mathcal P) = J_1 + J_2 = \sum_{k=1}^{K_1} \big(\|\Delta\bm v^{(1)}_{k,1}\| + \|\Delta\bm v^{(1)}_{k,2}\|\big) + \sum_{k=1}^{K_2} \big(\|\Delta\bm v^{(2)}_{k,1}\| + \|\Delta\bm v^{(2)}_{k,2}\|\big) .
  \label{eq:q3obj}
\end{equation}

\subsubsection{约束条件分析}

两颗观测星的时间约束、脉冲约束和观测约束与第 6.1.3 节相同，且两星的累计速度增量预算相互独立，分别满足 $J_1 \le 300$~m/s 与 $J_2 \le 300$~m/s。两星的观测目标互不重复，即 $\pi^{(1)}$ 与 $\pi^{(2)}$ 的元素两两不同且交集为空。

对于 R2 的初始轨道状态 $\bm x_0^{(2)}$，题面要求 R2 为提前入轨、偏心率小于 0.001 的高轨卫星，并未给出半长轴的具体数值范围。本文依据目标星座的实际分布引入半长轴偏离上限参数 $\delta a$：由附件数据，30 颗蓝方目标的半长轴相对地球同步轨道半长轴 $\aGEO = 42164.17$~km 的偏离量介于 $-5.90$~km（B30）与 $+4.73$~km（B21）之间，红方 R1 的偏离量为 $+1.77$~km。同步带上蓝方其他高轨卫星和红方高轨卫星的半长轴偏离量同样较小，R2 作为红方提前入轨的高轨卫星，其初始轨道应与上述卫星同处同步带内。本文因此取 $\delta a = 5.90$~km（即目标半长轴的最大偏离量），并将偏心率与半长轴约束施加于 R2 在 $t_0$ 时刻的初始轨道：
\begin{equation}
  \mathcal X_{\mathrm{init}}(\delta a) = \Big\{\bm x_0^{(2)} \in \mathbb R^6 \;\Big|\; \|\bm e(\bm x_0^{(2)})\| < 0.001,\ \big|a(\bm x_0^{(2)}) - \aGEO\big| \le \delta a \Big\} ,
  \label{eq:q3init}
\end{equation}
其中 $a(\cdot)$ 与 $\bm e(\cdot)$ 分别是由位置速度矢量计算半长轴与偏心率矢量的解析公式。$t_0$ 之后的机动照常计入 R2 的速度增量，机动后的轨道只受机动约束与预算限制。为检验 $\delta a$ 取值的影响，后文将在 $\delta a = 3 \sim 19$~km 范围内给出最优总速度增量曲线 $J^*(\delta a)$。

\subsubsection{模型汇总}

沿用式~\eqref{eq:q2set} 的时间与脉冲可行集记号 $\mathcal Z_{K_1}$、$\mathcal Z_{K_2}$，问题三的优化模型汇总为
\begin{equation}
  \operatorname*{lexmax}_{\substack{\pi^{(1)},\,\pi^{(2)},\,\bm x_0^{(2)}\in\mathcal X_{\mathrm{init}}(\delta a)\\ \bm z^{(1)}\in\mathcal Z_{K_1},\,\bm z^{(2)}\in\mathcal Z_{K_2}}}\ \big(K_1 + K_2,\ -(J_1 + J_2)\big)
  \label{eq:q3model}
\end{equation}
\begin{equation*}
  \text{s.t.}\quad
  \begin{aligned}[t]
    & \bm\rho^{(i)}_{\pi^{(i)}_k}(s^{(i)}_k + j;\,\mathcal P) \in \mathcal C_{\pi^{(i)}_k}(s^{(i)}_k + j),\quad j = 0,\dots,60,\ k = 1,\dots,K_i,\ i \in \{1, 2\}, \\
    & \pi^{(1)} \cap \pi^{(2)} = \varnothing,\quad \pi^{(i)}_k \ne \pi^{(i)}_l\ (k \ne l) .
  \end{aligned}
\end{equation*}

\subsection{优化模型求解}

模型~\eqref{eq:q3model} 在问题二的基础上增加了目标分配和 R2 初始轨道设计两层决策。求解按三个递进步骤展开：首先按相对经度将 30 颗目标划分为东西两组以消除折返；其次为 R2 选定初始轨道初值，并在两星各自的任务集上应用三维 CW 物理状态动态规划与独立排程；最后建立本文主模型 \textbf{M3-D}，将两星的全部脉冲矢量、脉冲施加时刻以及 R2 的初始六维轨道状态一并纳入序列凸优化做连续联合精修。

\subsubsection{按经度划分东西双星任务集}

由图~\ref{fig:targets} 可知，30 颗目标在同步带上分布在 R1 初始经度的东西两侧：西侧有 10 颗（B1、B2、B3、B5、B8、B9、B17、B18、B26、B30），跨度约 $18.8^\circ$；东侧有 20 颗（B4、B6、B7、B10 至 B16、B19 至 B25、B27 至 B29），跨度约 $18.9^\circ$。以 R1 初始经度为界，将西侧 10 颗分配给 R1，东侧 20 颗分配给 R2，两颗星的任务随即解耦。

这样分工有两个直接好处。第一，R1 从西侧组的东端出发，只需单向向西漂移；R2 可以直接布设在东侧组的西端，单向向东漂移，两颗星都不需要折返。第二，每颗星要覆盖的相位跨度都从 $38^\circ$ 减半到约 $19^\circ$，而可用时间仍是 30 天，平均漂移率大幅降低。由式~\eqref{eq:wmax}，低漂移率使卫星能够平缓接近那些交线偏离公共拱线较远的目标，避免了问题二中频繁急减速再加速的高额代价。

\subsubsection{R2 初始轨道初设与三维 CW 物理状态排程}

在初始轨道的粗搜索阶段，先将 R2 置于与 R1 共面的近圆轨道上，对半长轴偏离 $\Delta a$、倾角 $i$、升交点赤经 $\Omega$ 和初始相位 $\varphi_0$ 做网格扫描。每个候选初始轨道上，利用第 6.2.1 节的漂移相位束搜索与逐段精修方法，分别求出 R1 观测西侧 10 颗目标、R2 观测东侧 20 颗目标的初始方案。

这一阶段（记作模型 M3-A）在新数据上得到 R1 累计 18.94~m/s、R2 累计 16.91~m/s、合计 35.85~m/s 的全覆盖初始方案。随后对两星各自应用第 6.2.2 节的全计划序列凸优化（记作模型 M3-B），两星分别降至 18.15~m/s 和 15.96~m/s，合计 34.10~m/s。我们也对两星各自运行了 1~h 的并行大邻域搜索以检验访问顺序：R1 的顺序没有找到改进；R2 搜出的新顺序经两阶段联合优化后为 16.67~m/s，高于原序，说明经度分组后两星各自的访问顺序已接近局部最优。

\subsubsection{脉冲时刻—矢量—初始轨道全耦合序列凸优化（M3-D）}

上述基线模型（M3-C）虽已放开 R2 初始六维状态 $\bm x_0^{(2)}$ 做固定时刻序列凸优化（合计 33.46~m/s），但仍有两处显著改进空间：一是 R1 在 M3-A/M3-B 中仅用前 18.6 天便草草收尾，未能充分利用后 11 天的时间窗口来降低漂移率；二是两星所有脉冲的施加时刻仍被冻结在初始排程上。

为此，本文将问题二的 M2-D 框架完整推广至双星系统，提出主模型 \textbf{M3-D}：（1）在 R1 西侧 10 颗目标上运行零先验三维 CW 物理状态动态规划（3D-CW-VND-SE-DP），重新分配 30 天内的最佳交会日期并放开脉冲时刻与矢量联合优化 $(\Delta\bm v_m, t_m)$；（2）在 R2 东侧 20 颗目标上，将初始六维状态修正 $\delta\bm x_0^{(2)} \in \mathbb R^6$、全部 40 次脉冲矢量修正 $\delta\Delta\bm v_m$ 及脉冲时刻修正 $\delta t_m$ 同步纳入二阶锥规划子问题：
\begin{equation}
  \bm r_{\mathrm{R2}}(t) \approx \bar{\bm r}_{\mathrm{R2}}(t) + \bm\Phi_{rx}(t, 0)\,\delta\bm x_0^{(2)} + \sum_{t_m < t} \Big(\bm\Phi_{rv}(t, t_m)\,\delta\Delta\bm v_m + \bm\Phi_{rt}(t, t_m)\,\delta t_m\Big) ,
  \label{eq:q3lin}
\end{equation}
其中 $\bm\Phi_{rx}$、$\bm\Phi_{rv}$ 与 $\bm\Phi_{rt}$ 均由精确二体传播 24 核并行有限差分得到。初始轨道的偏心率与半长轴约束在当前初始状态 $\bar{\bm x}_0^{(2)}$ 处线性化为
\begin{equation}
  \big\|\bm e(\bar{\bm x}_0^{(2)}) + \bm J_e\,\delta\bm x_0^{(2)}\big\| \le e_{\max},\qquad
  \aGEO - \delta a \le a(\bar{\bm x}_0^{(2)}) + \bm J_a\,\delta\bm x_0^{(2)} \le \aGEO + \delta a ,
  \label{eq:q3initlin}
\end{equation}
其中 $\bm J_e \in \mathbb R^{3\times 6}$ 和 $\bm J_a \in \mathbb R^{1\times 6}$ 分别是偏心率矢量和半长轴关于初始状态的雅可比矩阵。连同式~\eqref{eq:socp} 的观测凸集、单次脉冲上界和相邻脉冲最小间隔约束 $(t_{m+1}+\delta t_{m+1}) - (t_m+\delta t_m) \ge 3601$~s，每轮子问题仍是标准二阶锥规划，由 Clarabel 高效求解。对两星分别运行“时刻—矢量（—初始轨道）联合 VC-SCP $\to$ 固定时刻精确精修”闭环，每轮均用精确二体模型在 0.05~s 网格上严格认证观测约束，并用解析公式精确检验 $e(\bm x_0^{(2)}) < 0.001$ 与 $|a(\bm x_0^{(2)}) - \aGEO| \le \delta a$。

\subsection{结果的分析与验证}

在主方案 $\delta a = 5.90$~km 下，M3-D 仅用约 6~min 即求得 30 天全覆盖双星协同最优方案：两星累计速度增量合计降至 \textbf{28.11~m/s}，其中 R1 观测西侧 10 颗目标消耗 \textbf{13.50~m/s}（占单星预算 4.5\%，较 M3-C 下降 25.6\%），R2 观测东侧 20 颗目标消耗 \textbf{14.61~m/s}（占单星预算 4.9\%，较 M3-C 下降 4.6\%）。与问题二 M2-D 单星方案的 159.33~m/s 相比，双星协同使总速度增量大幅下降了 \textbf{82.4\%}。

优化得到的 R2 初始轨道状态与经典根数见表~\ref{tab:q3r2init}。R2 初始半长轴为 42158.275~km，较 $\aGEO$ 低 5.895~km，紧贴 $\delta a = 5.90$~km 的下界以自带向东自然漂移；偏心率为 0.000482，远小于 0.001 的上限；倾角和升交点赤经与 R1 保持一致，便于利用公共拱线附近的自然交会。

\begin{table}[!htb]
  \centering
  \caption{问题三 R2 初始轨道设计结果（M3-D，$\delta a = 5.90$~km，$t_0$ 时刻 ICRF 系）}
  \label{tab:q3r2init}
  \zihao{5}
  \begin{tabular}{lc@{\hspace{24pt}}lc}
    \toprule
    ICRF 状态分量 & 数值 & 经典轨道根数 & 数值 \\
    \midrule
    $x_0$ / km & $-30222.705$ & 半长轴 $a$ / km & 42158.275 \\
    $y_0$ / km & $-29342.917$ & 相对偏离 $a - \aGEO$ / km & $-5.895$ \\
    $z_0$ / km & $1748.022$ & 偏心率 $e$ & 0.000482 \\
    $v_{x0}$ / (km/s) & $2.130479$ & 倾角 $i$ / ($^\circ$) & 4.3795 \\
    $v_{y0}$ / (km/s) & $-2.208222$ & 升交点赤经 $\Omega$ / ($^\circ$) & 76.963 \\
    $v_{z0}$ / (km/s) & $-0.197109$ & 相对 R1 初始相位差 & $+4.07^\circ$（东侧） \\
    \bottomrule
  \end{tabular}
\end{table}

两星各自的观测序列和单次观测速度增量见表~\ref{tab:q3plan}。在 M3-D 中，R1 充分利用全部 30 天时间窗（第 29.57 天结束），以极平缓的向西漂移完成 10 颗目标，共 20 次脉冲（有效脉冲仅 12 次），单次最大脉冲 4.96~m/s，最小脉冲间隔高达 35770.8~s（近 10~h），且对 B30 几乎无需专门机动；R2 在第 29.52 天完成东侧 20 颗目标，共 40 次脉冲（有效脉冲 22 次），单次最大脉冲仅 2.83~m/s，对 B4、B7、B12、B29 均几乎实现零机动自然飞越。

\begin{table}[!htb]
  \centering
  \caption{问题三双星观测序列（M3-D 主方案；左：R1 西侧 10 颗，右：R2 东侧 20 颗）}
  \label{tab:q3plan}
  \zihao{5}\setlength{\tabcolsep}{3pt}
  \begin{tabular}{clcc@{\hspace{10pt}}clcc}
    \toprule
    序号 & 目标 & 观测开始（UTC） & $\Delta v$/(m/s) & 序号 & 目标 & 观测开始（UTC） & $\Delta v$/(m/s) \\
    \midrule
    1 & B2 & 01-02 13:43:57 & 5.84 & 1 & B4 & 01-01 14:07:49 & 0.00 \\
    2 & B1 & 01-03 14:22:11 & 0.55 & 2 & B11 & 01-04 13:51:05 & 1.16 \\
    3 & B3 & 01-07 14:16:47 & 0.15 & 3 & B7 & 01-05 14:58:30 & 0.00 \\
    4 & B5 & 01-08 14:15:18 & 0.45 & 4 & B6 & 01-06 17:24:14 & 0.42 \\
    5 & B9 & 01-12 13:52:10 & 0.37 & 5 & B10 & 01-08 13:25:40 & 0.05 \\
    6 & B8 & 01-16 13:54:09 & 1.82 & 6 & B12 & 01-11 13:11:50 & 0.00 \\
    7 & B17 & 01-21 13:45:35 & 1.96 & 7 & B13 & 01-12 13:07:24 & 0.78 \\
    8 & B18 & 01-26 11:35:03 & 1.01 & 8 & B15 & 01-14 12:38:48 & 1.37 \\
    9 & B30 & 01-27 17:46:02 & 0.00 & 9 & B14 & 01-15 12:32:34 & 0.40 \\
    10 & B26 & 01-30 13:42:20 & 1.36 & 10 & B16 & 01-15 17:21:55 & 0.26 \\
    & & & & 11 & B19 & 01-18 12:20:35 & 1.62 \\
    & & & & 12 & B20 & 01-19 12:26:16 & 1.98 \\
    & & & & 13 & B21 & 01-20 17:13:59 & 1.22 \\
    & & & & 14 & B23 & 01-22 12:06:58 & 0.70 \\
    & & & & 15 & B22 & 01-23 17:57:51 & 0.05 \\
    & & & & 16 & B24 & 01-25 11:28:07 & 0.14 \\
    & & & & 17 & B25 & 01-26 15:36:10 & 1.62 \\
    & & & & 18 & B27 & 01-28 11:29:06 & 0.03 \\
    & & & & 19 & B29 & 01-29 13:04:07 & 0.00 \\
    & & & & 20 & B28 & 01-30 12:21:18 & 2.83 \\
    \midrule
    \multicolumn{4}{l}{R1 合计 13.50~m/s（10 颗）} & \multicolumn{4}{l}{R2 合计 14.61~m/s（20 颗）；双星总计 28.11~m/s} \\
    \bottomrule
  \end{tabular}
\end{table}

双星的巡游路线和各次观测速度增量如图~\ref{fig:q3route} 所示。与图~\ref{fig:q2route} 的单星路线相比，R1 和 R2 从队列中部分头向西、向东匀速推进，轨迹近似两条斜率更平缓的直线，全程均铺满 30 天且没有折返。全部 30 颗目标单次观测速度增量的中位数从问题二的 2.79~m/s 骤降至 0.50~m/s；除 R1 建立初始向西漂移的首颗目标 B2（5.84~m/s）外，其余 29 颗目标的观测代价全部低于 2.9~m/s。

\begin{figure}[!htb]
  \centering
  \begin{subfigure}[b]{0.49\textwidth}
    \centering
    \includegraphics[width=\linewidth]{fig12a_q3_route.pdf}
    \caption{双星相对经度随时间的变化}
  \end{subfigure}\hfill
  \begin{subfigure}[b]{0.49\textwidth}
    \centering
    \includegraphics[width=\linewidth]{fig12b_q3_dv.pdf}
    \caption{双星各次观测的速度增量}
  \end{subfigure}
  \caption{问题三的双星协同巡游方案}
  \label{fig:q3route}
\end{figure}

表~\ref{tab:q3ablation} 对比了问题三四个递进模型的结果。在经度分组与离散网格初设（M3-A，35.85~m/s）的基础上，两星固定时刻序列凸优化（M3-B）降至 34.10~m/s；放开 R2 初始轨道做连续联合优化（M3-C）降至 33.46~m/s；进一步引入三维 CW 动态规划与脉冲时刻—矢量—初始轨道全耦合优化（M3-D），R1 降至 13.50~m/s、R2 降至 14.61~m/s，双星总计 \textbf{28.11~m/s}，较 M3-A 下降 \textbf{21.6\%}，较 M3-C 再降 \textbf{16.0\%}。

\begin{table}[!htb]
  \centering
  \caption{问题三各模型与求解阶段的对比}
  \label{tab:q3ablation}
  \zihao{-5}\setlength{\tabcolsep}{3pt}
  \begin{tabular}{lcccc}
    \toprule
    模型阶段 & R1/(m/s) & R2/(m/s) & 合计/(m/s) & $\delta a$/km \\
    \midrule
    M3-A：经度分组 + 网格扫描 + 逐段求解 & 18.94 & 16.91 & 35.85 & $\ge 7.80$ \\
    M3-B：M3-A + 两星固定时刻 SCP & 18.15 & 15.96 & 34.10 & $\ge 7.81$ \\
    M3-C：R2 初始轨道联合优化（固定时刻） & 18.15 & 15.31 & 33.46 & $5.90$ \\
    \textbf{M3-D（主方案）：3D-CW-DP + 联合 VC-SCP} & \textbf{13.50} & \textbf{14.61} & \textbf{28.11} & $\mathbf{5.90}$ \\
    \bottomrule
  \end{tabular}
\end{table}

为检验本文引入的参数 $\delta a$ 对结论的影响，我们在 M3-C 基线上于 $\delta a = 3 \sim 19$~km 范围内扫描求解了最优总速度增量 $J^*(\delta a)$，结果见图~\ref{fig:q3da} 与表~\ref{tab:q3da}。$J^*(\delta a)$ 随 $\delta a$ 单调不增：$\delta a \le 15$~km 时最优半长轴均贴在下界 $\aGEO - \delta a$ 上，说明半长轴约束起作用；当 $\delta a$ 增至约 19.01~km 后半长轴不再下降，$J^*(\delta a)$ 饱和于 33.27~m/s。从 $\delta a = 3$~km 到饱和点 19.01~km，两星总速度增量仅变化 0.29~m/s（0.9\%），$\delta a = 5.90$~km 与饱和点仅差 0.19~m/s（不足 0.6\%），说明本文的结论对 $\delta a$ 的取值不敏感。

\begin{figure}[!htb]
  \centering
  \includegraphics[width=0.68\textwidth]{fig13_q3_da.pdf}
  \caption{半长轴偏离上限 $\delta a$ 对双星最优总速度增量 $J^*(\delta a)$ 的影响（M3-C 基线扫描）}
  \label{fig:q3da}
\end{figure}

\begin{table}[!htb]
  \centering
  \caption{不同 $\delta a$ 下的双星最优解与 R2 初始轨道参数（M3-C 扫描）}
  \label{tab:q3da}
  \zihao{5}\setlength{\tabcolsep}{4pt}
  \begin{tabular}{ccccc}
    \toprule
    $\delta a$ / km & R2 / (m/s) & 两星合计 / (m/s) & 初始 $a - \aGEO$ / km & 初始偏心率 $e$ \\
    \midrule
    3.00 & 15.418 & 33.566 & $-2.99$ & 0.00051 \\
    5.00 & 15.345 & 33.493 & $-4.99$ & 0.00047 \\
    \textbf{5.90（基准）} & \textbf{15.312} & \textbf{33.460} & $-5.89$ & 0.00045 \\
    8.00 & 15.236 & 33.383 & $-7.99$ & 0.00041 \\
    15.00 & 15.150 & 33.298 & $-14.99$ & 0.00070 \\
    $\ge 19.01$（饱和点） & 15.126 & 33.274 & $-19.01$ & 0.00089 \\
    \bottomrule
  \end{tabular}
\end{table}

M3-D 主方案的约束检验与 ATK 独立复核结果见表~\ref{tab:q3check}。两星均按 0.05~s 网格做解析重放认证（ERFA 几何/视太阳双口径），并在 ATK 4.2.0 中通过分段接力数值积分独立复核，R1 10/10、R2 20/20 全部满足约束。全程 R1 与 R2 的最低地心距分别为 42116.1~km 和 42046.4~km，始终处于地球同步轨道带附近。

\begin{table}[!htb]
  \centering
  \caption{问题三双星主方案约束检验（M3-D）}
  \label{tab:q3check}
  \zihao{5}\setlength{\tabcolsep}{4pt}
  \begin{tabular}{lccc}
    \toprule
    检验项 & 要求 & R1（西侧 10 颗） & R2（东侧 20 颗） \\
    \midrule
    观测目标数 & 合计 30 颗不重复 & 10 / 10 & 20 / 20 \\
    累计速度增量 & 单星 $\le 300$~m/s & 13.50~m/s & 14.61~m/s \\
    单次速度增量 & $\le 30$~m/s & 最大 4.96~m/s & 最大 2.83~m/s \\
    相邻脉冲间隔 & $\ge 3600$~s & 最小 35770.8~s & 最小 3950.9~s \\
    每次观测前机动次数 & $\le 2$ & 最多 2 & 最多 2 \\
    R2 初始 $e$ / $|a-\aGEO|$ & $< 0.001$ / $\le 5.90$~km & — & 0.00048 / 5.90~km \\
    观测距离 & $10 \sim 50$~km & $10.05 \sim 49.95$~km & $10.04 \sim 49.95$~km \\
    最大光照角 / 侦测角 & $\le 45^\circ$ & $44.95^\circ$ / $44.95^\circ$ & $44.95^\circ$ / $44.95^\circ$ \\
    ATK 独立数值积分复核 & 全部满足 & 10 / 10 & 20 / 20 \\
    \bottomrule
  \end{tabular}
\end{table}
